\begin{summarybox}
\textbf{\textcolor{CSummary}{一分钟记牢}}

\textbf{\Large 先看轴在不在，再比端点谁更远。}

\begin{itemize}[leftmargin=*,itemsep=0.33em]
 \item \textbf{开口向上：}最低先找轴（轴不在区间就找近端）；\textbf{最高一定在远端}。
 \item \textbf{开口向下：}最高先找轴（轴不在区间就找近端）；\textbf{最低一定在远端}。
 \item \textbf{端点碰轴：}$m=h$ 或 $n=h$。这控制顶点值能否取到。
 \item \textbf{中点碰轴：}$\dfrac{m+n}{2}=h$。这控制两端点函数值何时相等、较远端何时换边。
\end{itemize}

\textbf{\textcolor{CSummary}{看到题目这样下手}}

\begin{center}
\begin{tikzpicture}[>=stealth, every node/.style={align=center,draw=teal!55!black,rounded corners=3pt,fill=teal!6,inner sep=5pt,font=\small}]
 \node (one) at (0,0) {求轴 $x=h$\\看开口};
 \node (two) at (3.35,0) {区间是否\\包含 $h$？};
 \node (three) at (6.75,0) {比较两端\\到轴的距离};
 \draw[->,thick,teal!70!black] (one)--(two);
 \draw[->,thick,teal!70!black] (two)--(three);
\end{tikzpicture}

\end{center}

\textbf{\textcolor{CSummary}{最常用的数字例子}}

$f(x)=(x-2)^2$、区间 $[t,t+2]$：\textbf{最小值}按 $t=0,2$ 分段；\textbf{最大值}按 $t=1$ 分段。请记住这\textbf{两类不同的分界点}，不要只会画“轴左、跨轴、轴右”三种图。

\textbf{\textcolor{CSummary}{使用前检查}}

函数为固定二次函数 $a\neq0$，区间是包含两端点的闭区间。遇到边界等号，代入相邻公式核对；开区间不能直接照搬最值结论。
\end{summarybox}
